\documentclass[11pt]{amsart}
\usepackage{calc,amssymb,amsmath,ulem,stmaryrd,fullpage}
\usepackage{enumerate}
\usepackage{alltt}
\RequirePackage[dvipsnames,usenames]{color}
\let\mffont=\sf
\normalem
%\input{kmacros3.sty}
\input{xy}
\xyoption{all}
\usepackage{hyperref}
\usepackage{graphicx}
\usepackage[all,cmtip]{xy}
\usepackage{verbatim}
\usepackage[left=0.95in,top=0.95in,right=0.95in,bottom=0.95in]{geometry}
\newcommand{\tauCohomology}{T}
\newcommand{\FCohomology}{S}


\theoremstyle{theorem}
%\newtheorem*{mainthm*}{Main Theorem}

\newcommand{\note}[1]{\marginpar{\sffamily\tiny #1}}

\def\todo#1{\textcolor{Mahogany}%
{\sffamily\footnotesize\newline{\color{Mahogany}\fbox{\parbox{\textwidth-15pt}{#1}}}\newline}}

\renewcommand{\O}{\mathcal O}
\newcommand{\ie}{{\itshape i.e.} }
\newcommand{\roundup}[1]{\lceil #1 \rceil}
\newcommand{\rounddown}[1]{\lfloor #1 \rfloor}
\renewcommand{\to}[1][]{\xrightarrow{\ #1\ }}
\newcommand{\onto}[1][]{\protect{\xrightarrow{\ #1\ }\hspace{-0.8em}\rightarrow}}
\newcommand{\into}[1][]{\lhook \joinrel \xrightarrow{\ #1\ }}
%\newcommand{DBDual}[1]{{\underline{\omega}_{#1}}}
\newcommand{\DBDual}[1]{{{\uuline \omega}{}^{\mydot}_{#1}}}
\newcommand{\Tt}{{\mathfrak{T}}}
%\newcommand{\bn}{{\mathfrak{n}} }
\newcommand{\sol}[1]{ \vskip 6pt{\color{blue} {\bf Solution:} #1} \vskip 6pt}

\newcommand{\bR}{{\mathbb{R}}}
\newcommand{\va}{{\vec{a}}}
\newcommand{\vb}{{\vec{b}}}
\newcommand{\vzero}{{\vec{0}}}
\newcommand{\vi}{{\vec{i}}}
\newcommand{\vj}{{\vec{j}}}
\newcommand{\vk}{{\vec{k}}}
\newcommand{\vh}{{\vec{h}}}
\newcommand{\vr}{{\vec{r}}}
\newcommand{\vu}{{\vec{u}}}
\newcommand{\vv}{{\vec{v}}}
\newcommand{\vw}{{\vec{w}}}
\newcommand{\vx}{{\vec{x}}}
\newcommand{\vy}{{\vec{y}}}
\newcommand{\vz}{{\vec{z}}}
\newcommand{\vT}{{\vec{T}}}
\newcommand{\vN}{{\vec{N}}}
\newcommand{\vF}{{\vec{F}}}
\newcommand{\vG}{{\vec{G}}}
\newcommand{\bF}{\mathbb{F}}
\newcommand{\bQ}{\mathbb{Q}}
\newcommand{\bZ}{\mathbb{Z}}
\newcommand{\bC}{\mathbb{C}}
\newcommand{\Gal}{\textnormal{Gal}}
\newcommand{\Aut}{\textnormal{Aut}}
\newcommand{\Emb}{\textnormal{Emb}}
\newcommand{\Tr}{\textnormal{Tr}}
\newcommand{\N}{\textnormal{N}}
\newcommand{\Hom}{\textnormal{Hom}}


\begin{document}
\title{HW \#4 -- Math 6320\\Spring 2015}
\author{Due: Thursday February 26th}
\maketitle

\begin{enumerate}
\item Determine the Galois group of $x^4-25$ over $\bQ$.
\item Let $K$ be a field of characteristic $\neq 2$.  Suppose that $\alpha, \beta \in K$.  Figure out exactly when $K(\sqrt{\alpha}) = K(\sqrt{\beta})$.  Use this to determine whether or not $\bQ(\sqrt{1 - \sqrt{2}}) = \bQ(i, \sqrt{2})$.
\item Let $K = \bQ(a^{1/n})$ where $a \in \bQ_{> 0}$ and that $x^n - a$ is irreducible so that $[K : \bQ] = n$.  Suppose that $E$ is a subfield of $K$ with $[E : \bQ] = d$.  Prove that $E = \bQ(a^{1/d})$.  
    \vskip 3pt
    \emph{Hint:} Consider $N_{K/E}(a^{1/n}) \in E$ (remember, $N_{K/E}(a^{1/n})$ was defined in the previous homework).
\item  Suppose that $m, n > 0$ are integers.  What is $(\bZ/m\bZ) \otimes (\bZ/n\bZ)$?  Of course your answer will depend on $m$ and $n$.
\item  Let $A$ be a ring, $M$ and $A$-module and $I$ an ideal.  Show that $M \otimes_A (A/I) \cong M/IM$.  
\item  Suppose that $R$ is a local ring and $M, N$ are finitely generated $R$-modules.  Prove that if $M \otimes_R N = 0$ then $M = 0$ or $N = 0$.  Find counter examples to this statement if $R$ is nonlocal or if $M, N$ are not finitely generated.
    \vskip 3pt
    \emph{Hint:} Use Nakayama's lemma.
\item Let $R$ be a ring $\neq 0$ with $R^{n} \cong R^{m}$ for some integers $m,n > 0$.  Show that $m = n$.
    \vskip 3pt
    \emph{Hint:}  Use Nakayama's lemma.
\item Suppose that $R$ is a ring and that $L, M, N$ are $R$-modules.  Show that $\Hom_R(L \otimes M, N) \cong \Hom_R(L, \Hom_R(M, N))$.  Additionally, choose one of the modules $L, M, N$ and show that this isomorphism is functorial in that variable.  In other words, if you choose $L$ and if $L \to L'$ is a module map, show that
    \[
    \xymatrix{
        \Hom_R(L' \otimes_R M, N) \ar[r] \ar@{<->}[d]^{\sim} & \Hom_R(L \otimes_R M, N) \ar@{<->}[d]^{\sim}  \\
        \Hom_R(L', \Hom_R(M, N)) \ar[r] & \Hom_R(L, \Hom_R(M, N))
    }
    \]
    commutes where the vertical maps are the isomorphisms from earlier in this problem and the horizontal maps are the ones coming from the contravariant nature of $\Hom$.  
    \vskip 3pt
    \emph{Hint:}  This is easier than you might think, try writing down where something has to go.
\end{enumerate}

\end{document}

